\section{形成性研究}

本文开展了一项形成性研究，以探索数据分析人员在图表创作过程中的挑战与需求。该研究已获得所在机构伦理审查委员会（Institutional Review Board, IRB）的批准，所有收集数据仅用于科学研究目的。参与者在研究开始前均已充分了解研究内容并签署知情同意，同时研究过程保证其数据严格保密，且仅用于本研究。

\subsection{参与者}

本文共招募了11名参与者，其中女性5名、男性6名，平均年龄为31.5岁。参与者按职业背景可分为三类：3名学生（P1-P3）、3名学者（E1、P4、P5）和5名行业从业者（E2、P6-P9）。参与者平均具有9.6年的数据可视化经验。其中，两名专家自愿参与形成性研究，其余参与者在完成研究后获得15美元报酬。参与者的详细信息如表\ref{tab:formative_participant} 所示。

\begin{table}[htbp]
\centering
\caption{形成性研究参与者信息}
\label{tab:formative_participant}
\small
\setlength{\tabcolsep}{3pt}
\begin{tabular}{lllll}
\toprule
参与者 & 身份 & 性别 & 使用场景 & 年龄 \\
\midrule
E1 & 教授 & 男 & 数据可视化、学术写作、数据分析 & 41 \\
E2 & 数据科学家 & 男 & 数据分析 & 43 \\
P1 & 本科生 & 女 & 课程作业、文档写作 & 21 \\
P2 & 硕士研究生 & 女 & 课程作业、文档写作、数据分析 & 24 \\
P3 & 硕士研究生 & 男 & 课程作业、文档写作、数据分析 & 26 \\
P4 & 博士研究生 & 女 & 学术写作、数据分析 & 28 \\
P5 & 博士研究生 & 男 & 学术写作、数据分析 & 31 \\
P6 & 数据叙事从业者 & 女 & 文档写作、数据分析 & 32 \\
P7 & 数据从业者 & 男 & 数据分析 & 33 \\
P8 & 新闻从业者 & 女 & 新闻写作、数据分析 & 33 \\
P9 & 数据从业者 & 男 & 数据分析 & 35 \\
\bottomrule
\end{tabular}
\end{table}

\subsection{研究程序}

本次形成性研究采用半结构化访谈方法，访谈通过线上与线下相结合的方式进行。线上访谈使用飞书或腾讯会议，线下访谈采用面对面交流。每位参与者均接受独立访谈，访谈目标是深入探讨图表推荐功能的必要性。

访谈内容围绕四个核心方面展开：参与者在图表构建方面的个人经验、典型图表创作流程、创作过程中遇到的挑战，以及其评估图表有效性的方法。每次访谈时长为35至60分钟。线上访谈通过视频录制和文字纪要进行记录，线下访谈则通过书面笔记与现场录音结合的方式进行记录。

\subsection{研究发现}

通过收集与分析半结构化访谈数据，本文对参与者反馈进行了系统整理与主题归纳，并提炼出若干关键发现。

\subsubsection{F1：图表选择过载}

访谈结果显示，参与者在选择最合适的数据可视化图表时，经常会产生犹豫或困惑。这种不确定性主要来自可选图表类型过多，使其难以判断哪一种图表最适合当前任务（P1、P4、P6、P8）。虽然参与者熟悉柱状图、折线图和饼图等基础图表，但面对大量选项时，他们往往默认选择这些熟悉类型（P8）。当数据复杂度提高时，用户的不确定性进一步增加，不清楚哪种图表能够更好地呈现数据关系。即使他们知道箱线图、漏斗图等较高级的图表类型，也常因不确定其适用性而避免使用这些图表（P6、P7）。因此，用户往往通过多个基础图表来简化复杂数据，这会导致信息被拆散，难以揭示数据中的深层关系和洞察（P7）。

P6是一位数据叙事从业者，她提到：“我手里有不同渠道的用户来源、转化率和最终消费额数据，我想在一张图里展示这个完整流程以及每个环节的占比，但我完全不知道应该用什么图。最后只能做三个独立的柱状图，但看起来并不容易理解。”

记忆并理解多种图表类型的适用场景本身就是一项专业技能，这对非专家用户或初学者而言会带来较高的认知负担。用户也常担心使用不常见图表会增加与受众沟通的成本，甚至影响他人对其专业性的判断，因此倾向于采取更保守的选择。传统工具（如 Excel）通常将所有图表类型不加区分地展示给用户，缺少足够的引导与上下文信息，反而进一步增加了图表选择难度。

\subsubsection{F2：低效的试错循环}

非专家用户的图表构建过程并不是线性的、由意图驱动的设计流程，而往往表现为低效的试错循环。该过程的典型特征是，用户在数据操作的句法层面与可视化表达的语义层面之间反复切换，其核心机制造成了明显的效率损耗。

这一循环通常从用户基于自身对数据的语义理解和初步分析意图尝试生成图表开始。然而，由于用户心智模型与工具对数据结构的句法要求之间普遍存在差异，初次尝试经常导致视觉编码失败。随后，面对这一障碍，用户的认知焦点会发生显著转移：从“我希望数据表达什么”这一较高层次的语义目标，转向“我应该如何重构数据以满足工具规范”这一较低层次的程序性任务。

P7是一位日常工作中频繁处理数据、但并非可视化设计专家的从业者，他描述道：“我经常为了做一个图，把同一个数据表复制粘贴好几份，用不同格式来尝试，看哪个能生成我想要的图。感觉一半时间都花在这种繁琐工作上。最令人头疼的是，好不容易把图做出来了，老板说‘换个角度展示看看’，然后我又要重复刚才所有步骤。”

这种循环往复的过程不仅大量消耗用户时间，也带来了较重的认知负担，并可能使用户在持续挫败中降低甚至放弃最初的分析目标。因此，最终生成的可视化图表往往并不是数据洞察的最佳视觉表达，而是用户在探索空间中由于时间或认知资源耗尽而选择的“局部最优解”或“满意解”。这一现象揭示出现有可视化工具在引导用户完成从数据到视觉的顺畅映射方面仍存在不足。

\subsubsection{F3：从数据到视觉的转译挑战}

用户在数据可视化过程中面临的核心智力挑战，是完成从特定领域问题到抽象视觉编码的认知跨越。许多情况下，用户清楚自己想通过数据讲述什么故事，或者想回答什么业务问题，例如“如何展示各产品线在不同季度对总销售额的贡献占比”。但是，用户并不清楚如何将这个“故事”或“问题”转译为合适的数据可视化方案，例如“堆叠柱状图”。用户的思维停留在具体业务语境中，而可视化工具的语言是抽象图表类型，二者之间存在难以跨越的鸿沟。

P9是一位数据从业者，他提到：“我知道我想表达什么，我们 A 产品的销售额远超预期，B 产品则没有达标。但我不知道怎么在一张图里把‘实际’、‘目标’和‘差距’这三个信息都清楚展示出来。”

其根本原因在于，非专家用户普遍缺乏一种被称为“可视化素养”（Visualization Literacy）的多层次综合能力。这种素养的核心，是能够将具体业务问题抽象为通用分析任务，例如比较、构成和关系等，并掌握将这些任务映射到最佳视觉编码方案的图式知识。当前主流工具以“图表为中心”的设计范式加剧了这一困境：工具预设用户已经具备这种素养，仅提供执行层面的“图表选项板”，从而在最关键的认知转译环节忽略了非专家用户的核心需求。

\subsubsection{F4：简单化的评价标准}

非专家用户对可视化图表的评价并不主要基于分析层面的严谨性，而是更容易受到审美吸引力和信息传递效率这两个表层修辞目标驱动。前者关注图表是否美观、是否能在汇报或展示中呈现“专业感”；后者则关注观众是否能够快速理解并接受一个预设结论，也就是所谓的“一眼看懂”。因此，图形保真度（如坐标轴完整性）、信息密度以及是否可能引发误读等更深层次的分析质量，往往被忽略。

P8是一位新闻从业者，她说：“我做的图主要是给公众看的，所以最重要的是把结论清楚标出来，让他们一眼就能看到重点，其他信息越少越好。”

这一现象的根源可归结为用户角色目标与专业知识水平的共同作用。在大多数场景下，用户扮演的是“观点传达者”而非“数据探索者”的角色，其核心目标是以最高效率说服观众接受既定结论。这种“修辞”导向的目标，使用户自然会优先选择简洁、美观的设计；而“图形诚信”（Graphical Integrity）原则的普遍缺失，则使他们难以辨别有效简化与误导性失真之间的关键差异，甚至可能在无意识中将“可视化反模式”视为实现高效沟通的捷径。

\subsection{系统设计需求}

基于上述访谈发现，本文提出三项设计原则。这些原则旨在指导智能图表推荐系统的开发，并为用户从意图构建到生成准确、有效的视觉表达提供认知脚手架。

\noindent\textbf{意图先行原则。}
系统的主要交互模型应从询问用户“你想创建什么图表？”转向询问“你想回答什么问题？”（F1、F3）。用户输入包含数值信息的自然语言文本，系统解析其意图并推荐合适图表~\cite{latif2022kori,masry2023unichart,tian2024chartgpt,wang2022towards}。

$\bullet$ \textbf{R1：自然语言意图识别。}
系统必须具备自然语言处理能力，能够准确解析用户输入文本，并从文本中识别核心意图，例如比较数值、观察趋势、探索关系等。

$\bullet$ \textbf{R2：可解释的推荐。}
系统应将推荐结果完整展示给用户，包括被推荐图表的原始案例以及对应排名。

\noindent\textbf{无摩擦探索原则。}
系统设计应尽量降低用户在探索过程中的认知成本和操作成本，从而鼓励迭代优化，而不是让用户在挫败感中妥协（F2）~\cite{luo2021natural}。

$\bullet$ \textbf{R3：并行草稿预览。}
系统应改变串行试错流程，转而同时渲染多个候选视觉编码方案，并以低保真预览形式呈现，使用户能够快速进行横向比较与选择。

$\bullet$ \textbf{R4：探索历史版本比较。}
系统应自动记录探索路径，并以缩略图列表等可视化方式呈现。用户可以随时保存某个满意版本，在不同版本之间进行并排比较，或一键回溯到先前任一状态，从而将线性的“试错-撤销”过程转化为非线性的、可比较的探索过程。

\noindent\textbf{规范化初始状态。}
考虑到系统能力边界，系统输出优化的核心应聚焦于为用户提供一个高质量起点~\cite{stokes2022striking}。这能够在后续用户自定义过程中提供认知支持，而不是执行完全自动化的审查与美化（F4）。

$\bullet$ \textbf{R5：以最佳实践作为默认输出。}
系统基于推荐生成的初始图表必须严格遵循图形诚信与美学设计的最佳实践。例如，柱状图的 Y 轴默认从0开始，默认配色方案应对色盲用户友好，标签应清晰且无重叠。这确保用户图表创作的起点是高质量、无误导的标准模板。

\newpage
